% GENERATED FILE. Edit canonical lessons, then run npm run generate:books.
% canonical-source-hash: fnv1a64:36d2d732e2c018ae
% canonical-lessons: JA-C204-ryokou, JA-C204-ryokan, JA-C204-kenbutsu, JA-C204-yoru, JA-C204-yado

\chapter{Trips, Inns and Sights}
\label{ch:ja-trips-inns-and-sights}
\hlchaptermodality{204}

\begin{chapteropening}
\textbf{By the end of this chapter:} \emph{I can talk about a trip and where I will stay.}

Complete the last lesson of chapter 204: I can talk about a trip and where I will stay.
\end{chapteropening}

\section[\texorpdfstring{\ja{りょこう} (ryokō)}{りょこう (ryokō)}]{\ja{りょこう} (ryokō) — a trip; travel}

\label{lesson:JA-C204-ryokou}



\begin{quote}
Before the new one: say the Japanese for a seat, then the Japanese for to give up.
\end{quote}

\subsection*{You'll want to know: \ja{りょこう}}
\textbf{\ja{りょこう}} — \emph{ryokō} — \textquotedblleft{}a trip; travel\textquotedblright{}.

\begin{quote}
\textbf{\ja{らいげつ} \ja{りょこうに} \ja{いきます。}} — \emph{raigetsu ryokō ni ikimasu} — \textquotedblleft{}I am going on a trip next month.\textquotedblright{}
\end{quote}

\begin{cousinweb}
\textbf{\ja{りょこう}} (\emph{ryokō}) is Sino-Japanese (kango): two Chinese characters read with their on'yomi, readings borrowed from Chinese. They are \emph{ryo}, \textquotedblleft{}journey\textquotedblright{} (Middle Chinese \emph{ljoX}) and \emph{kō}, \textquotedblleft{}to go\textquotedblright{} (\emph{haeng}) (KANJIDIC2; Middle Chinese in Baxter's transcription).
\end{cousinweb}

\begin{grammarlens}[title={What you've built}]
One more word for everyday life at A2.
\end{grammarlens}

\subsection*{Your turn}
\begin{itemize}
\raggedright
  \item \emph{Say it:} \emph{ryokō}
  \item \emph{Say it:} \emph{ryokō}, once more
  \item \emph{Say it:} \emph{raigetsu ryokō ni ikimasu} — I am going on a trip next month.
  \item \emph{Recall:} say the Japanese for a seat, then the Japanese for to give up, then say \emph{ryokō} again
\end{itemize}

\begin{tcolorbox}[breakable,colback=teal!4,colframe=teal!35!black,title={Before you move on}]
What does \ja{りょこう} mean? (\textquotedblleft{}A trip; travel\textquotedblright{}.) Say it once more.
\end{tcolorbox}
\section[\texorpdfstring{\ja{りょかん} (ryokan)}{りょかん (ryokan)}]{\ja{りょかん} (ryokan) — a Japanese inn}

\label{lesson:JA-C204-ryokan}



\begin{quote}
Before the new one: say the Japanese for to give up, then the Japanese for a trip.
\end{quote}

\subsection*{You'll want to know: \ja{りょかん}}
\textbf{\ja{りょかん}} — \emph{ryokan} — \textquotedblleft{}a Japanese inn\textquotedblright{}.

\begin{quote}
\textbf{\ja{りょかんで} \ja{ゆっくり} \ja{やすみます。}} — \emph{ryokan de yukkuri yasumimasu} — \textquotedblleft{}I rest well at the inn.\textquotedblright{}
\end{quote}

\begin{cousinweb}
\textbf{\ja{りょかん}} (\emph{ryokan}) is Sino-Japanese (kango): two Chinese characters read with their on'yomi, readings borrowed from Chinese. They are \emph{ryo}, \textquotedblleft{}journey\textquotedblright{} (Middle Chinese \emph{ljoX}) and \emph{kan}, \textquotedblleft{}building, lodge\textquotedblright{} (\emph{kwanH}) (KANJIDIC2; Middle Chinese in Baxter's transcription).
\end{cousinweb}

\begin{grammarlens}[title={What you've built}]
One more word for everyday life at A2.
\end{grammarlens}

\subsection*{Your turn}
\begin{itemize}
\raggedright
  \item \emph{Say it:} \emph{ryokan}
  \item \emph{Say it:} \emph{ryokan}, once more
  \item \emph{Say it:} \emph{ryokan de yukkuri yasumimasu} — I rest well at the inn.
  \item \emph{Recall:} say the Japanese for to give up, then the Japanese for a trip, then say \emph{ryokan} again
\end{itemize}

\begin{tcolorbox}[breakable,colback=teal!4,colframe=teal!35!black,title={Before you move on}]
What does \ja{りょかん} mean? (\textquotedblleft{}A Japanese inn\textquotedblright{}.) Say it once more.
\end{tcolorbox}
\section[\texorpdfstring{\ja{けんぶつ} (kenbutsu)}{けんぶつ (kenbutsu)}]{\ja{けんぶつ} (kenbutsu) — sightseeing}

\label{lesson:JA-C204-kenbutsu}



\begin{quote}
Before the new one: say the Japanese for a trip, then the Japanese for a Japanese inn.
\end{quote}

\subsection*{You'll want to know: \ja{けんぶつ}}
\textbf{\ja{けんぶつ}} — \emph{kenbutsu} — \textquotedblleft{}sightseeing\textquotedblright{}.

\begin{quote}
\textbf{\ja{ともだちと} \ja{まちの} \ja{けんぶつに} \ja{いきます。}} — \emph{tomodachi to machi no kenbutsu ni ikimasu} — \textquotedblleft{}I go sightseeing in town with a friend.\textquotedblright{}
\end{quote}

\begin{cousinweb}
\textbf{\ja{けんぶつ}} (\emph{kenbutsu}) is Sino-Japanese (kango): two Chinese characters read with their on'yomi, readings borrowed from Chinese. They are \emph{ken}, \textquotedblleft{}to see\textquotedblright{} (Middle Chinese \emph{kenH}) and \emph{butsu}, \textquotedblleft{}thing\textquotedblright{} (\emph{mjut}) (KANJIDIC2; Middle Chinese in Baxter's transcription).
\end{cousinweb}

\begin{grammarlens}[title={What you've built}]
One more word for everyday life at A2.
\end{grammarlens}

\subsection*{Your turn}
\begin{itemize}
\raggedright
  \item \emph{Say it:} \emph{kenbutsu}
  \item \emph{Say it:} \emph{kenbutsu}, once more
  \item \emph{Say it:} \emph{tomodachi to machi no kenbutsu ni ikimasu} — I go sightseeing in town with a friend.
  \item \emph{Recall:} say the Japanese for a trip, then the Japanese for a Japanese inn, then say \emph{kenbutsu} again
\end{itemize}

\begin{tcolorbox}[breakable,colback=teal!4,colframe=teal!35!black,title={Before you move on}]
What does \ja{けんぶつ} mean? (\textquotedblleft{}Sightseeing\textquotedblright{}.) Say it once more.
\end{tcolorbox}
\section[\texorpdfstring{\ja{よる} (yoru)}{よる (yoru)}]{\ja{よる} (yoru) — to drop in; to stop by}

\label{lesson:JA-C204-yoru}



\begin{quote}
Before the new one: say the Japanese for a Japanese inn, then the Japanese for sightseeing.
\end{quote}

\subsection*{You'll want to know: \ja{よる}}
\textbf{\ja{よる}} — \emph{yoru} — \textquotedblleft{}to drop in; to stop by\textquotedblright{}.

\begin{quote}
\textbf{\ja{えきの} \ja{ちかくの} \ja{ほんやに} \ja{よります。}} — \emph{eki no chikaku no hon'ya ni yorimasu} — \textquotedblleft{}I drop in at the bookshop near the station.\textquotedblright{}
\end{quote}

Polite \textbf{\ja{よります}}, \emph{yorimasu}; a godan verb.

\begin{cousinweb}
\textbf{\ja{よる}} (\emph{yoru}) is a native Japanese word (wago), inherited rather than borrowed from Chinese; its kanji was attached later as a kun'yomi, and its older origin is not settled (JMdict, KANJIDIC2).
\end{cousinweb}

\begin{grammarlens}[title={What you've built}]
One more word for everyday life at A2.
\end{grammarlens}

\subsection*{Your turn}
\begin{itemize}
\raggedright
  \item \emph{Say it:} \emph{yoru}
  \item \emph{Say it:} \emph{yoru}, once more
  \item \emph{Say it:} \emph{eki no chikaku no hon'ya ni yorimasu} — I drop in at the bookshop near the station.
  \item \emph{Recall:} say the Japanese for a Japanese inn, then the Japanese for sightseeing, then say \emph{yoru} again
\end{itemize}

\begin{tcolorbox}[breakable,colback=teal!4,colframe=teal!35!black,title={Before you move on}]
What does \ja{よる} mean? (\textquotedblleft{}To drop in; to stop by\textquotedblright{}.) Say it once more.
\end{tcolorbox}
\section[\texorpdfstring{\ja{やど} (yado)}{やど (yado)}]{\ja{やど} (yado) — lodging; a place to stay}

\label{lesson:JA-C204-yado}



\begin{quote}
Before the new one: say the Japanese for sightseeing, then the Japanese for to drop in.
\end{quote}

\subsection*{You'll want to know: \ja{やど}}
\textbf{\ja{やど}} — \emph{yado} — \textquotedblleft{}lodging; a place to stay\textquotedblright{}.

\begin{quote}
\textbf{\ja{やどは} \ja{うみの} \ja{ちかくです。}} — \emph{yado wa umi no chikaku desu} — \textquotedblleft{}Our lodging is near the sea.\textquotedblright{}
\end{quote}

\begin{cousinweb}
\textbf{\ja{やど}} (\emph{yado}) is a native Japanese word (wago), inherited rather than borrowed from Chinese; its kanji was attached later as a kun'yomi, and its older origin is not settled (JMdict, KANJIDIC2).
\end{cousinweb}

\begin{grammarlens}[title={What you've built}]
One more word for everyday life at A2.
\end{grammarlens}

\subsection*{Your turn}
\begin{itemize}
\raggedright
  \item \emph{Say it:} \emph{yado}
  \item \emph{Say it:} \emph{yado}, once more
  \item \emph{Say it:} \emph{yado wa umi no chikaku desu} — Our lodging is near the sea.
  \item \emph{Recall:} say the Japanese for sightseeing, then the Japanese for to drop in, then say \emph{yado} again
\end{itemize}

\begin{tcolorbox}[breakable,colback=teal!4,colframe=teal!35!black,title={Before you move on}]
What does \ja{やど} mean? (\textquotedblleft{}Lodging; a place to stay\textquotedblright{}.) Say it once more.
\end{tcolorbox}
